\chapter{Memory Systems and Cognitive Persistence}

Claude Code's memory systems implement DCF's ``language as infrastructure'' principle:

\section{CLAUDE.md (Project Memory)}

\begin{lstlisting}
# CLAUDE.md - Project: MyApp

## Project Overview
[Brief description for context]

## Architecture Decisions
- Using PostgreSQL for X reason
- Chose React over Vue because...

## Conventions
- Test files: *.test.ts
- Use conventional commits

## DCF Preferences
- Always present plans before implementing
- Challenge my assumptions
- Show trade-offs explicitly
\end{lstlisting}

\section{User Memories (Personal Cognitive Profile)}

Accumulated understanding of how you work:
\begin{itemize}
    \item Preferred coding style
    \item Common project types
    \item Review depth preferences
\end{itemize}

\section{Serena Memories (Code Intelligence)}

Project-specific technical understanding:
\begin{itemize}
    \item Architecture patterns discovered during onboarding
    \item Symbol relationships
    \item Suggested commands
\end{itemize}

\section{Session Continuity: Preparing for Compaction}\label{sec:session-continuity}\index{Session Continuity}\index{Compaction Strategy}

Long sessions eventually hit context limits. When Claude Code compacts the conversation, it generates a summary to preserve continuity---but the quality of that summary depends on what's available to summarize. \textbf{Proactive documentation before compaction} dramatically improves session continuity.

\subsection{The Problem}

Context compaction is lossy. Details that seemed obvious during the session may not survive summarization:
\begin{itemize}
    \item Decisions made and their rationale
    \item Files modified and why
    \item Approaches considered and rejected
    \item Next steps that were identified but not yet started
\end{itemize}

Without explicit capture, you may resume a compacted session and find yourself re-discovering context you already had.

\subsection{The Practice: Pre-Compaction Documentation}

When you sense a session is getting long or complex, create a working document that captures session state. This serves two purposes: (1) it creates explicit content for the compaction summary to reference, and (2) it forces you to synthesize what actually happened.

\textbf{Recommended structure:}

\begin{lstlisting}
# Session Findings

## Completed Work
- What was accomplished
- Which files were modified and why
- Decisions made with rationale

## Open Questions
- Unresolved issues
- Clarifications still needed

## Recommended Next Steps
- Prioritized list of remaining work
- Dependencies between tasks

## Context That Matters
- Assumptions being made
- Constraints discovered during session
- Anything non-obvious that future-you needs to know
\end{lstlisting}

\subsection{When to Capture}

\begin{itemize}
    \item \textbf{Before natural breaks}: Lunch, end of day, context switch
    \item \textbf{After completing a logical chunk}: Feature done, review complete
    \item \textbf{When the session feels ``heavy''}: Many files touched, complex decisions made
    \item \textbf{Before requesting compaction}: If you're about to run \texttt{/compact}
\end{itemize}

\subsection{Gitignore Strategy}

Session documents are working artifacts, not permanent documentation. Consider gitignoring them:

\begin{lstlisting}
# .gitignore
SESSION_NOTES.md
SESSION_FINDINGS.md
\end{lstlisting}

This keeps your repository clean while maintaining session continuity. The content either gets incorporated into permanent artifacts (CLAUDE.md, code comments, documentation) or becomes obsolete.

\subsection{Connection to DCF Principles}

Pre-compaction documentation is \textbf{anticipatory calibration} applied to session management. You're forming an explicit model of ``what matters about this session'' before the compaction forces you to rely on an automated summary. The act of writing clarifies your own understanding---the thinking mirror principle at work.

It's also \textbf{scaffolding that fades}: the session document supports continuity during active work, then becomes unnecessary once the work is complete and captured in permanent artifacts.

\section{Shared Context Infrastructure}\label{sec:shared-context}\index{Shared Context Infrastructure}\index{Context Engineering Repository}

Individual memory systems (CLAUDE.md, personal memories) extend naturally to \textbf{team-level shared context}. Organizations that maximize AI benefits create infrastructure for collective context---preventing the siloing of AI knowledge.

\subsection{The Problem: AI Knowledge Silos}

Without deliberate infrastructure, AI benefits fragment:
\begin{itemize}
    \item Each team member develops their own effective prompts
    \item Successful patterns aren't shared
    \item Failed experiments are repeated across the team
    \item No institutional learning about what works
\end{itemize}

This mirrors traditional knowledge silos, amplified by AI's rapid iteration cycles.

\subsection{The Solution: Context Engineering Repositories}

A \textbf{context engineering repository} serves as a centralized place for team AI knowledge:

\begin{lstlisting}
/context-engineering/
  /prompts/
    - code-review-prompt.md
    - architecture-analysis.md
    - debugging-approach.md
  /CLAUDE.md-templates/
    - project-type-a.md
    - project-type-b.md
  /learnings/
    - what-works.md
    - what-doesnt-work.md
  /experiments/
    - current-experiments.md
    - results.md
\end{lstlisting}

\subsection{Goals of Shared Context Infrastructure}

\begin{enumerate}
    \item \textbf{Team alignment}: Shared vocabulary and approaches
    \item \textbf{Iterative improvement}: Prompts evolve through collective experience
    \item \textbf{Transparency}: ``This is what we're trying, here's what we're learning''
    \item \textbf{Collaboration}: Others can try, adapt, and improve shared patterns
    \item \textbf{Prevent siloing}: Useful processes become organizational assets, not individual secrets
    \item \textbf{Knowledge centralization}: One place to look for ``how do we use AI for X''
\end{enumerate}

\subsection{Implementation Patterns}

\textbf{Pull request workflow for prompts:} Treat prompt improvements like code changes. Propose, review, merge. This creates accountability and learning records.

\textbf{Regular sync meetings:} Brief (15 min) sessions to share ``what I learned about AI this week.'' Surfaces patterns that wouldn't otherwise propagate.

\textbf{Experiments registry:} Track what's being tried, by whom, with what results. Prevents duplicate effort and builds institutional knowledge.

\begin{quotebox}
\textbf{Key insight:} The most effective organizations treat AI collaboration as a \emph{team sport}, not an individual skill. Shared context infrastructure makes this possible.
\end{quotebox}

